Automated eye-blink extraction from electroencephalography (EEG) is useful both for artifact handling and for physiological analyses, but the best known dedicated workflow, \blinker{}, remains centered on MATLAB. We report \tool, a Python implementation intended for direct use within MNE-Python pipelines, and test whether it reproduces the MATLAB reference behavior at scale rather than only approximately. We compared Python and MATLAB event tables on two public EEG resources representing distinct acquisition contexts: 74 locally staged EDF/FIF-compatible recordings from the Murat 2018 motor-imagery dataset and all 34 available EDF recordings from the SEED-VLA/VRW fatigue-driving dataset. Comparisons used harmonized blink boundaries, a \num{20}-sample tolerance window, and strict and lenient event-level metrics. Full-dataset reruns exposed two non-trivial SEED boundary failures: a truncated terminal-edge blink that survived to the final event table, and a floating-point amplitude-gate mismatch near the recording end. After correcting these cases, while rejecting a broader intermediate fix that introduced a new regression, the final Python implementation achieved exact parity with MATLAB on all 108 recordings. Across the combined benchmark, both implementations reported \num{156871} blink events, with zero Python-only and zero MATLAB-only events, and all strict and lenient macro and micro precision, recall, F1 and accuracy metrics equal to 1.0. These results establish \tool{} as a faithful, Python-native route to \blinker-style ocular-index extraction for reproducible EEG analysis.
